%% Experimental Setup & Baselines
\section*{Experimental setup and baselines}

The baseline is repeated PrexSyn sampling. We draw $N$ additional candidates from the same specification and evaluate them identically, setting $N$ at or beyond each method's diversity plateau as determined by the pilot saturation study. We classify a variant as a hit if AiZynthFinder finds a synthetic route and the variant conserves properties on both axes below. We treat thresholds as parameters (Tanimoto $\geq$ 0.6, 0.7, and 0.85; desirability $\geq$ 0.8) so that conclusions do not depend on any single cutoff.

We assess property conservation on two axes. \textit{Substructural} conservation measures ECFP4 Tanimoto similarity between the variant and the specification fingerprint, capturing the fraction of shared circular substructure features. \textit{Physicochemical} conservation computes per-descriptor deltas (MW $\pm$50\,Da, CLogP $\pm$1.0, TPSA $\pm$20\,\AA$^2$, HBD $\pm$1, HBA $\pm$2, rotatable bonds $\pm$2), maps each to a 0--1 desirability score, and aggregates them by geometric mean. Both axes match PrexSyn's conditioning variables. We include ESPsim \cite{bib12}, which captures 3D shape and electrostatic similarity, as a supplementary diagnostic.

The primary metrics are hit rate and expansion factor. Hit rate is the fraction of synthesizable variants meeting conservation thresholds. Expansion factor divides unique method hits by unique baseline hits, directly measuring marginal value: a value $>$1 indicates viable candidates unreachable by repeated sampling alone. Secondary metrics include drug-likeness via QED \cite{bib13} and chemical diversity ($1 -$ mean pairwise ECFP4 Tanimoto among hits).

We stratify all metrics by baseline-quality bin ($<$0.5, 0.5--0.7, 0.7--0.85, 0.85--1.0). Pairwise Jaccard overlap between hit sets identifies complementary tool combinations.
